\textbf{Width sweep.} Table~\ref{tab:sweep} and Figure~\ref{fig:sweep} train uniform DR at seven widths. Nominal return is at the cap for every width (499.9 at $w=3$, the only value below the cap). Robustness rises steeply from $w=0$ (386.6) to $w=1$ (470.4), reaches 493.3 at $w=1.5$ and 493.2 at $w=2$, and is 490.0 at $w=3$ with larger seed variance (std 8.2). The only candidate for a ``knee'' is therefore a flat or very slightly decreasing robustness beyond $w\approx1.5$, a difference far smaller than the seed spread; nominal performance does not decrease at all. So within the range we tested, we found no width at which nominal performance dropped faster than robustness improved. Note \texttt{ret\_robust} for $w\ge1.5$ is partly in-range, so the plateau also reflects the evaluation ceiling.

\begin{table}[t]\centering\caption{Sweep over fixed training width $w$ (mean $\pm$ std over 5 seeds).}\label{tab:sweep}\resizebox{\columnwidth}{!}{\begin{tabular}{lcccc}
\toprule
Train width $w$ & ret\_nominal & ret\_w100 & ret\_ood & ret\_robust \\
\midrule
0 & 500.0 $\pm$ 0.0 & 389.1 $\pm$ 28.8 & 243.2 $\pm$ 21.9 & 386.6 $\pm$ 16.6 \\
0.25 & 500.0 $\pm$ 0.0 & 440.2 $\pm$ 29.2 & 292.4 $\pm$ 11.3 & 418.1 $\pm$ 11.0 \\
0.5 & 500.0 $\pm$ 0.0 & 479.6 $\pm$ 15.4 & 344.0 $\pm$ 37.2 & 444.5 $\pm$ 14.9 \\
1 & 500.0 $\pm$ 0.0 & 498.5 $\pm$ 1.7 & 412.0 $\pm$ 34.3 & 470.4 $\pm$ 11.7 \\
1.5 & 500.0 $\pm$ 0.0 & 500.0 $\pm$ 0.0 & 479.9 $\pm$ 3.3 & 493.3 $\pm$ 1.1 \\
2 & 500.0 $\pm$ 0.0 & 500.0 $\pm$ 0.0 & 479.7 $\pm$ 5.1 & 493.2 $\pm$ 1.7 \\
3 & 499.9 $\pm$ 0.2 & 498.6 $\pm$ 3.1 & 470.9 $\pm$ 23.0 & 490.0 $\pm$ 8.2 \\
\bottomrule
\end{tabular}
}\end{table}

\begin{figure}[t]\centering\includegraphics[width=0.48\columnwidth]{sweep_ret_nominal.pdf}\includegraphics[width=0.48\columnwidth]{sweep_ret_robust.pdf}\caption{Nominal return (left) and \texttt{ret\_robust} (right) versus training width.}\label{fig:sweep}\end{figure}

\textbf{Curriculum gate.} Table~\ref{tab:abl} varies the success threshold. The four variants are indistinguishable at 5 seeds: \texttt{ret\_robust} differs by under 8 units, within one seed standard deviation, and no paired comparison against gate 0.8 is below $p=0.18$ (gate 0.5: 0.19, gate 0.95: 0.22, no gate: 0.85); the ordering is also non-monotone in the threshold, so we offer no explanation for it. All variants reach the cap width. The gate is thus not shown to help in this task, where even narrow ranges are solved quickly; its intended benefit (avoiding ranges too hard to learn) is not exercised. Figure~\ref{fig:curr} shows the width trajectory of the main curriculum, which reaches the cap about halfway through training.

\begin{table}[t]\centering\caption{Curriculum gate ablation (mean $\pm$ std over 5 seeds; main method uses gate 0.8).}\label{tab:abl}\resizebox{\columnwidth}{!}{\begin{tabular}{lcccc}
\toprule
Method & ret\_nominal $\uparrow$ & ret\_ood $\uparrow$ & ret\_robust $\uparrow$ & $n$ \\
\midrule
Curriculum, no gate & \textbf{500.00 $\pm$ 0.00} & \textbf{414.97 $\pm$ 21.75} & \textbf{471.62 $\pm$ 7.30} & 5 \\
Curriculum, gate 0.95 & 500.00 $\pm$ 0.00 & 391.20 $\pm$ 23.70 & 463.51 $\pm$ 8.04 & 5 \\
Curriculum, gate 0.8 & 500.00 $\pm$ 0.00 & \underline{412.89 $\pm$ 38.93} & \underline{470.84 $\pm$ 13.07} & 5 \\
Curriculum, gate 0.5 & \underline{500.00 $\pm$ 0.00} & 394.57 $\pm$ 21.19 & 464.69 $\pm$ 7.13 & 5 \\
\bottomrule
\end{tabular}
}\end{table}

\begin{figure}[t]\centering\includegraphics[width=0.85\columnwidth]{curriculum_width.pdf}\caption{Range half-width $w$ during training for the success-gated curriculum, mean and spread over 5 seeds.}\label{fig:curr}\end{figure}
